\uuid{BiPN}
\exo7id{3444}
\titre{exo7 3444}
\auteur{quercia}
\organisation{exo7}
\datecreate{2010-03-10}
\isIndication{false}
\isCorrection{false}
\chapitre{Déterminant, système linéaire}
\sousChapitre{Applications}
\cpge{mpsi-mp2i.espaces-vectoriels.c, pcsi.espaces-vectoriels.c, ptsi.espaces-vectoriels.c, mpsi-mp2i.determinants.b, pcsi.determinants.a, ptsi.determinants.c, mpsi-mp2i.determinants.d, pcsi.calcul-matriciel-systemes-lineaires.d, ptsi.calcul-matriciel-systemes-lineaires.d, bcpst1.espace-vectoriel-kn.a, pt.determinants.a, bcpst2.espaces-vectoriels.a}
\module{Algèbre}
\niveau{L2}
\difficulte{}

\contenu{
\question{
Soient $E$ un ensemble quelconque et
${f_1, \dots, f_n} : E \to K$ des fonctions.

Montrer par récurrence sur $n$ que la famille $(f_1, \dots, f_n)$ est libre
dans $ K^E$ si et seulement s'il existe des éléments $x_1, \dots, x_n$ de $E$
tels que $\det\bigl( f_i(x_j) \bigr) \ne 0$.
}
}
